\chapter*{پیشگفتار}
\addcontentsline{toc}{chapter}{پیشگفتار}
\markboth{پیشگفتار}{پیشگفتار}

برنامه‌نویسی یعنی گفت‌وگو با رایانه. شما چیزی می‌خواهید، آن را با زبانی
که رایانه می‌فهمد می‌نویسید و رایانه، بی‌چون‌وچرا و با سرعتی باورنکردنی،
همان را انجام می‌دهد. سال‌هاست که این گفت‌وگو تنها به زبان انگلیسی ممکن
بوده است؛ واژه‌هایی مانند \code{if} و \code{while} و \code{print} برای
میلیون‌ها نفر نخستین سد راه یادگیری بوده‌اند، نه به این خاطر که
برنامه‌نویسی دشوار است، بلکه به این خاطر که پیش از رسیدن به اندیشه‌ی
برنامه‌نویسی، باید از دیوار یک زبان بیگانه بالا می‌رفتند.

زبان برنامه‌نویسی سلام برای برداشتن همین دیوار ساخته شده است. در سلام
می‌نویسید \fcode{اگر}، می‌نویسید \fcode{تکرار}، می‌نویسید \fcode{چاپ}؛ نام
متغیرها و تابع‌هایتان را به فارسی می‌گذارید و برنامه‌ای که می‌نویسید مانند
یک متن فارسی خوانده می‌شود. سلام یک زبان واقعی و کامل است، نه یک اسباب‌بازی
آموزشی: با آن می‌توان برنامه‌های بزرگ نوشت، صفحه‌ها و وب‌سایت‌ها ساخت و
حتی خود مترجم زبان سلام نیز با زبان سلام نوشته شده است. سلام املای انگلیسی
و عربی هم دارد و برنامه‌ای که به فارسی می‌نویسید، همان قدرت و همان سرعت را
دارد که اگر به انگلیسی می‌نوشتید. اما این کتاب درباره‌ی سلام فارسی است؛ از
این‌جا به بعد، همه چیز به زبان خودمان خواهد بود.

\section*{این کتاب برای چه کسی است؟}

برای هر کسی که می‌خواهد برنامه‌نویسی را از صفر بیاموزد. هیچ پیش‌نیازی لازم
نیست: نه آشنایی قبلی با برنامه‌نویسی، نه ریاضیات پیشرفته و نه زبان
انگلیسی. تنها چیزی که لازم دارید یک رایانه یا حتی یک گوشی با مرورگر است و
کمی حوصله. دانش‌آموزان دبیرستان، دانشجویان سال اول، معلمان، و هر بزرگسال
کنجکاوی که همیشه دلش می‌خواسته بداند «برنامه‌نویسی یعنی چه» مخاطب این کتاب
هستند. اگر پیش‌تر با زبان دیگری برنامه نوشته‌اید هم می‌توانید این کتاب را
بخوانید، اما احتمالاً بخش‌های نخست را سریع‌تر پشت سر خواهید گذاشت.

\section*{چیزی روی رایانه نصب نکنید}

یکی از تصمیم‌های مهم ما در نوشتن این کتاب این بود که خواننده هیچ‌گاه مجبور
نشود با «خط فرمان»، «ترمینال» یا نصب نرم‌افزار درگیر شود. همه‌ی برنامه‌های
کتاب را می‌توانید در ویرایشگر برخط سلام بنویسید و اجرا کنید:
\begin{center}
\url{https://salamlang.github.io/Salam/}
\end{center}
این ویرایشگر در مرورگر شما اجرا می‌شود، به هیچ سرویس دیگری وابسته نیست و
همان چیزی است که در فصل نخست با آن آشنا می‌شوید. کافی است نشانی را باز
کنید، زبان فارسی را برگزینید و بنویسید.

\section*{در این کتاب چه چیزهایی می‌آموزید؟}

مسیر این کتاب همان مسیر آشنای همه‌ی کتاب‌های خوب برنامه‌نویسی است، با این
تفاوت که در هر گام کد فارسی می‌نویسید:
\begin{itemize}
  \item در فصل \ref{ch:first} نخستین برنامه‌ی خود را می‌نویسید و یاد
        می‌گیرید چگونه چیزی را روی صفحه چاپ کنید.
  \item فصل‌های \ref{ch:variables} تا \ref{ch:operators} به متغیر، حافظه،
        انواع داده و محاسبه می‌پردازند؛ یعنی مصالح ساختمانی هر برنامه.
  \item در فصل \ref{ch:conditions} برنامه‌ها یاد می‌گیرند تصمیم بگیرند و
        در فصل \ref{ch:loops} یاد می‌گیرند کاری را بارها تکرار کنند.
  \item فصل \ref{ch:functions} به شما نشان می‌دهد چگونه برنامه را به
        قطعه‌های کوچک و قابل استفاده‌ی دوباره تقسیم کنید.
  \item فصل‌های \ref{ch:arrays} و \ref{ch:strings} درباره‌ی نگه‌داری
        فهرستی از داده‌ها و کار با متن هستند.
  \item و فصل \ref{ch:projects} همه‌ی این‌ها را در چند برنامه‌ی کوچک و
        کاربردی کنار هم می‌گذارد.
\end{itemize}
ما عمداً از الگوریتم‌های مشهور و برنامه‌های طولانی پرهیز کرده‌ایم. هدف این
کتاب این است که اصول پایه را چنان خوب یاد بگیرید که بعد از آن هر چیزی را
بتوانید خودتان بیاموزید. مثال‌ها همه ساده و ملموس‌اند: محاسبه‌ی معدل، صورت
حساب خرید، جدول تبدیل دما، شمارش معکوس و چیزهایی از این دست.

\section*{چگونه این کتاب را بخوانیم؟}

برنامه‌نویسی با خواندن یاد گرفته نمی‌شود؛ با نوشتن یاد گرفته می‌شود. از شما
می‌خواهیم هر برنامه‌ای را که در کتاب می‌بینید، خودتان در ویرایشگر تایپ کنید
(نه اینکه کپی کنید)، اجرایش کنید و خروجی را با آنچه در کتاب آمده مقایسه
کنید. سپس آن را دست‌کاری کنید: عددی را عوض کنید، خطی را حذف کنید، چیزی
اضافه کنید و ببینید چه می‌شود. بهترین معلم شما همین آزمایش‌کردن است.

در سراسر کتاب چند نوع جعبه می‌بینید. جعبه‌های \textbf{نکته} چیزی را روشن‌تر
می‌کنند، جعبه‌های \textbf{مراقب باشید} از اشتباه‌های رایج خبر می‌دهند،
جعبه‌های \textbf{مثال} یک کاربرد کوچک را نشان می‌دهند و جعبه‌های
\textbf{بیشتر بدانید} برای خوانندگان کنجکاو نوشته شده‌اند و می‌توانید در
خواندن نخست از آن‌ها بگذرید. پایان هر فصل تمرین‌هایی دارد؛ پاسخ شماری از
آن‌ها را در پیوست پایانی کتاب آورده‌ایم، اما خواهش می‌کنیم پیش از نگاه
کردن به پاسخ، خودتان تلاش کنید.

\section*{سخن نویسندگان}

این کتاب حاصل همکاری خالق زبان سلام و یک مدرس دانشگاه است؛ یکی از درون
زبان به آن نگاه کرده و دیگری از نگاه کلاس درس و دانشجویی که برای نخستین بار
با برنامه‌نویسی روبه‌رو می‌شود. امیدواریم این دو نگاه در کنار هم کتابی
ساخته باشد که هم درست باشد و هم مهربان. اگر در خواندن کتاب به اشکالی
برخوردید یا پیشنهادی داشتید، خوشحال می‌شویم آن را از راه صفحه‌ی پروژه‌ی
سلام در گیت‌هاب با ما در میان بگذارید.

\begin{flushleft}
سیدعلی محمدیه\\
مهدی اسماعیلی
\end{flushleft}
\clearpage
